\section{Correctness of the interval certificate}
\label{int:certificate}
\label{int:appendix}

This appendix explains the integer interval certificate for the strict
positivity of the polynomial~\eqref{wei:polynomial}.  The coefficients are
those in~\eqref{wei:coefficients}, with the matching parameter choices
in~\eqref{wei:parameters}.  The certified domain consists of every probability
vector on three states, every $\xi\in[3/100,1]$, and each of the three closed
$a$ intervals
\[
 [0,190/1000],\qquad [190/1000,448/1000],\qquad
 [448/1000,477/1000].
\]
The third interval uses the Poisson values $H=1/12$, $U=1/2$; the first
two use $H=1/16$, $U=15/32$.  The source is supplied unchanged as
\texttt{verification/interval.cpp}; see also
Section~\ref{cert:interval}.  Assertions and signed 128-bit integer
arithmetic are part of the certificate.  We first explain what its
arithmetic computes and then why normal completion proves the claimed
inequality over the entire domain.

\subsection{Coordinates, formulas, and boundary extensions}
\label{int:parametrization}

The quantities being checked are invariant under permutation of the
three coordinates.  After sorting a probability vector in decreasing
order, it has the representation
\begin{equation}
 p=\frac{(1,q^{18},r^{18})}{Z},\qquad
 Z=1+q^{18}+r^{18},\qquad 0\le r\le q\le1.
 \label{int:coordinates}
\end{equation}
Indeed, the largest probability is positive; take the eighteenth roots
of the other two probabilities divided by it.  The formulas below are
valid on the larger square $0\le q,r\le1$, which will allow rectangular
intervals that cross the diagonal $r=q$.

Set $m_i=3p_i$ and $v_i=m_i-1$.  Then
\[
 v=\left(\frac3Z-1,\frac{3q^{18}}Z-1,
                  \frac{3r^{18}}Z-1\right),\qquad
 x=\frac16\sum_{i=1}^3v_i^2,\qquad y=\frac12v_1v_2v_3.
\]
In particular, $-1\le v_i\le2$ everywhere on the square.  For positive
$q,r$, write $L_q=\log q$ and $L_r=\log r$.  The centered logarithms
$u_i=\log m_i-\frac13\sum_j\log m_j$ are
\[
 (-6(L_q+L_r),\ 12L_q-6L_r,\ 12L_r-6L_q).
\]
Consequently the functions used in the weighted argument are exactly
\begin{align}
 G&=\frac{3q^6r^6}{Z},& T&=1-G,\notag\\
 W&=\frac{216}{Z}q^6r^6(L_q^2+L_r^2-L_qL_r),\notag\\
 Y&=\frac{648}{Z}q^6r^6
    (2L_q^3+2L_r^3-3L_q^2L_r-3L_qL_r^2).
 \label{int:weighted-functions}
\end{align}
These expressions explain the constants $216$ and $648$ in
\texttt{coeff}.  The entries returned by \texttt{qt} are enclosures of
$q^{18}$ and $q^6(\log q)^j$, $0\le j\le3$, together with their
derivatives; the analogous entries are used for $r$.

For the edge functions put
\[
 e_i=1-av_i,\qquad G_e=(e_1e_2e_3)^{1/3},\qquad
 \mathcal L(t)=\begin{cases}-\log(1-t)/t,&t\ne0,\\1,&t=0.
 \end{cases}
\]
Define
\[
 b_i=-v_i\mathcal L(av_i),\qquad
 w_i=b_i-\frac13\sum_jb_j.
\]
For $a>0$, $b_i=\log(e_i)/a$, so the centered edge logarithms are
$aw_i$.  Since $\sum_i v_i=0$,
\[
 G_e^3=1-3a^2x-2a^3y.
\]
Factoring $1-G_e^3$ gives the formulas used in the code:
\begin{equation}
 \widehat T=\frac{3x+2ay}{1+G_e+G_e^2},\qquad
 \widehat W=\frac{G_e}{3}\sum_iw_i^2,\qquad
 \widehat Y=aG_ew_1w_2w_3.
 \label{int:edge-functions}
\end{equation}
Here $\widehat T=T_e/a^2$, $\widehat W=W_e/a^2$, and
$\widehat Y=Y_e/a^2$ when $a>0$.  The last identity uses
$\frac13\sum_iw_i^3=w_1w_2w_3$ for three centered coordinates.
Thus the code variables \texttt{Te}, \texttt{We}, and \texttt{Ye} are
the hatted quantities.  At $a=0$, the formulas give continuously
\[
 G_e=1,\qquad \widehat T=x,\qquad
 \widehat W=2x,\qquad \widehat Y=0.
\]
There is no division by $a$ in the implemented expressions.  Substituting
\eqref{int:weighted-functions} and~\eqref{int:edge-functions} into the
four entries returned by \texttt{coeff} gives, term by term,
$(D_0,D_1,D_2,D_3)$ in~\eqref{wei:coefficients}.  In the parameter
structure, \texttt{del}, \texttt{ch}, and \texttt{cu} mean
$\delta,H,U$, and the code forms $s_0=c_0+1/3$.

All boundary probability vectors are included.  Each function
$q^6(\log q)^j$, $0\le j\le3$, is assigned value zero at $q=0$.
Its derivative for $q>0$ is
\begin{equation}
 6q^5(\log q)^j+jq^5(\log q)^{j-1},
 \label{int:log-product-derivative}
\end{equation}
where the second term is omitted when $j=0$.  Both the function and
its derivative tend to zero at zero.  Thus the expressions have
continuous first derivatives up to the boundary; when $qr=0$ they give
$T=1$ and $W=Y=0$.  Since $Z\ge1$ and $e_i>0$ throughout the domain,
all four $D_i$ have the same regularity, including at $a=0$.

\subsection{Integer intervals and machine safety}
\label{int:arithmetic}

Let $S=2^{40}$.  An object \texttt{I} with integer endpoints $(l,h)$
represents $[l/S,h/S]$.  The functions \texttt{unit} and
\texttt{integer} represent a grid point and a small integer,
respectively.  With positive denominator, \texttt{rat}$(b,c)$ represents
\[
 \left[\frac{\lfloor Sb/c\rfloor}{S},
       \frac{\lceil Sb/c\rceil}{S}\right].
\]
The helpers \texttt{mf}$(b,c)$ and \texttt{floor\_div}$(b,c)$ compute
the mathematical floors of $bc/S$ and $b/c$, also for negative
numerators: the negative branch uses the negative of the corresponding
positive ceiling.  Upper endpoints are obtained by negating a floor of
a negated expression.  Addition and subtraction are exact on the
grid.  Multiplication chooses the appropriate endpoint extrema in
each sign case; when both intervals cross zero it takes the minimum
of the two potentially negative products and the maximum of the two
potentially positive products.  The square operation assigns lower
endpoint zero when its input crosses zero.

For a positive interval $[l/S,h/S]$, the reciprocal is enclosed by
\[
 \left[\frac{\lfloor S^2/h\rfloor}{S},
       \frac{\lceil S^2/l\rceil}{S}\right].
\]
Negative reciprocals use odd symmetry.  An assertion excludes intervals
containing zero.  Consequently all elementary interval operations
contain their exact real counterparts.  Intersecting a value interval
with a separately established range is also sound.  The two uses of
this operation, \texttt{clip}, are the range $[-1,2]$ for $v_i$ and
the logarithmic-product bounds proved below.

These statements concern exact integer operations, so absence of
overflow must be established before using their outputs.  Every
potentially growing endpoint calculation is passed through
\texttt{safe}, which asserts
\[
 -2^{62}<b<2^{62}
\]
before converting the signed 128-bit integer $b$ to a signed 64-bit
endpoint.  Inductively, two input endpoints have product of magnitude
less than $2^{124}$.  The product is formed after promoting to signed
128 bits, and the rounding offset $S-1$ still leaves its magnitude
below $2^{127}$.  Endpoint addition is also promoted before checking.
Reciprocal numerators are $S^2=2^{80}$; adding the denominator minus
one for upward rounding is safe.  Rational-construction numerators
and denominators in the prescribed runs are small fixed integers.
No unchecked negation or absolute value can encounter the minimum
signed 64-bit integer, because endpoints have magnitude below $2^{62}$.

The remaining raw integer expressions have separate small bounds.
All box coordinates lie in $[0,S]$ in grid units, so midpoint sums
are at most $2S$.  Cube-root search candidates are below $2S$,
their endpoint sums are at most $4S$, and the two cube-comparison
operands are below $8S^3=2^{123}$.  The fixed multiples of $S$, loop
indices, and initial-grid numerators fit signed 64 bits.
The unchecked visit counter is treated separately by the concrete
sanitized executions recorded below; the depth bound alone does not
keep its total number of increments within the signed 64-bit range.
Logarithm scaling uses at most forty doublings on the actual
arguments.  Thus a range assertion is never being asked to detect an
overflow that has already occurred: oversized derivative or value
enclosures stop at \texttt{safe} before narrowing.  All floating-point
casts in this source occur in failure diagnostics, not in acceptance
tests.

\subsection{Logarithms, divided logarithms, and cube roots}
\label{int:logs}

The function \texttt{atanh2} evaluates the fifteen terms
\[
 2\sum_{n=0}^{14}\frac{w^{2n+1}}{2n+1}
\]
by Horner's rule with outward interval operations.  For $|w|\le u<1$
the absolute remainder is at most
\[
 \frac{2u^{31}}{31(1-u^2)}.
\]
The code asserts $|w|\le1/3+4/S<3/8$.  With $u=3/8$, the remainder
bound, in grid units, is the exact rational number
\begin{equation}
 S\frac{2(3/8)^{31}}{31(1-(3/8)^2)}
 =\frac{617673396283947}{119978708822917120}
 <\frac1{100}<2.
 \label{int:atanh-tail}
\end{equation}
The padding $[-2/S,2/S]$ therefore covers the omitted tail; all
polynomial-evaluation errors are already covered by interval rounding.

To enclose $\log(x/S)$ for a positive integer endpoint $x$,
\texttt{ln} scales an interval containing $x/S$ by powers of two.
If $m$ is the resulting interval and $k$ the accumulated exponent,
then $x/S\in2^km$.  Multiplication by two is exact on this grid;
halving, when used, rounds outwards.  The identities
\[
 \log t=2\operatorname{atanh}\frac{t-1}{t+1},\qquad
 \log2=2\operatorname{atanh}\frac13
\]
give the enclosure after adding $k\log2$.  The domain assertion in
\texttt{atanh2} checks the needed range even after rounding.  All
positive arguments actually passed to \texttt{ln} are either coordinate
endpoints in $[1/S,1]$, or endpoints of $1-av_i$ bounded away from
zero, so the scaling loops terminate within the stated bounds.

The divided logarithm has the integral and power-series representations
\begin{equation}
 \mathcal L(t)=\int_0^1\frac{ds}{1-st}
             =\sum_{n=0}^{\infty}\frac{t^n}{n+1}
 \quad (|t|<1).
 \label{int:L-representation}
\end{equation}
For $|t|\le1/4$, \texttt{L\_endpoint} evaluates
$\sum_{n=0}^{22}t^n/(n+1)$ and its derivative
$\sum_{n=1}^{22}nt^{n-1}/(n+1)$ independently.  Bounding both
coefficient sequences by one bounds their absolute tails by
\begin{equation}
 \frac{(1/4)^{23}}{1-1/4}=\frac1{48S},\qquad
 \frac{(1/4)^{22}}{1-1/4}=\frac1{12S},
 \label{int:L-tails}
\end{equation}
respectively.  Each receives padding $[-1/S,1/S]$.  Outside this
small interval the exact identities
\[
 \mathcal L(t)=-\frac{\log(1-t)}t,\qquad
 \mathcal L'(t)=\frac{(1-t)^{-1}-\mathcal L(t)}t
\]
use only nonzero denominators.  On the entire real domain $t<1$,
differentiation of the integral representation gives
\[
 \mathcal L'(t)=\int_0^1\frac{s\,ds}{(1-st)^2}>0,\qquad
 \mathcal L''(t)=2\int_0^1\frac{s^2\,ds}{(1-st)^3}>0.
\]
Thus both $\mathcal L$ and $\mathcal L'$ increase, including at negative
arguments.  The lower enclosure at the left endpoint and upper
enclosure at the right endpoint suffice to enclose each function
on an interval.  This is precisely the rule used by \texttt{L}.

For a nonnegative integer $x<8S$, \texttt{cb\_floor} uses binary
search and the exact comparison
\[
 m^3\le xS^2
\]
to return $\lfloor S(x/S)^{1/3}\rfloor$.  Its initial search bracket
$[0,2S]$ is valid by $x<8S$.  Applying this function to the two input
endpoints, then increasing the upper result by one grid unit, encloses
the cube root.  The derivative enclosure is obtained from
$1/(3t^{2/3})$ using the already enclosed positive root.

The outward rounding of an $a$ endpoint enlarges the prescribed bands
by less than $1/S$, and never introduces negative $a$.  Hence throughout
the actual boxes $a\le477/1000+1/S$.  Even the interval enclosures of
$1-av_i$, after multiplication rounding and clipping $v_i$ to $[-1,2]$,
have lower endpoint greater than $45/1000$ and upper endpoint less
than $1478/1000$.  The successive product of the three factors has
lower endpoint at least $(45/1000)^3-3/S>0$ and upper endpoint at most
$(1478/1000)^3+3/S<8$.  These bounds justify the logarithm,
divided-logarithm, reciprocal, and cube-root domains on every initial
box and descendant.  The cube-root lower enclosure is positive as
well, so its derivative reciprocal is valid.

\subsection{Logarithmic products and derivative propagation}
\label{int:derivatives}

The zero endpoints in~\eqref{int:coordinates} are treated without
evaluating their logarithms.  For $0<q\le1$, set $t=-\log q\ge0$.
For integers $0\le j\le3$ and real $h\ge5$,
\[
 q^{h/2}|\log q|^j=e^{-ht/2}t^j\le1.
\]
For $j=0$ this is immediate.  For $j>0$ its maximum is
$(2j/(he))^j\le1$, since $2j/h\le6/5<e$.
It follows that
\begin{equation}
 |q^5(\log q)^j|\le q^{5/2}\le q^2,\qquad
 |q^6(\log q)^j|\le q^3.
 \label{int:qt-bound}
\end{equation}
For even $j$ the products are nonnegative, and for odd $j$ they are
nonpositive.  These are the exact external ranges used by
\texttt{qt\_direct} to intersect the power-and-logarithm enclosures.
When the lower endpoint is zero, the same ranges directly supply
an enclosure, with no logarithm call.  When the lower endpoint is
positive, monotonicity of $\log$ supplies its interval, and intersecting
its upper bound with $0$ is valid because $q\le1$.  The derivative
returned for every logarithmic product is exactly the interval
version of~\eqref{int:log-product-derivative}.  The derivative of
$q^{18}$ is $18q^{17}$.  Cached enclosures depend only on the interval
endpoints; \texttt{qt} puts their derivatives in the appropriate
$q$ or $r$ gradient coordinate.

An object \texttt{J} consists of a value interval and three derivative
intervals, corresponding to $(q,r,a)$.  Constants have zero gradients
and variables have the appropriate unit gradient.  Addition,
multiplication, squares, and reciprocals apply the identities
\[
 \partial_i(fg)=f\,\partial_i g+g\,\partial_i f,\qquad
 \partial_i(f^2)=2f\,\partial_i f,\qquad
 \partial_i(f^{-1})=-f^{-2}\partial_i f.
\]
Cube roots and $\mathcal L$ use the derivative enclosures just proved
and the chain rule.  Induction over these operations proves that the
value and every first partial derivative of the actual formula are
enclosed throughout the box.  Tightening a value enclosure by an
external range does not replace the underlying function by a clipped
function: its previously enclosed derivatives remain valid, and
the tighter values can be used in subsequent product and chain rules.
All these statements extend to zero coordinate endpoints by the
continuous derivative extensions already established.

\subsection{Bernstein coefficients and the acceptance test}
\label{int:bernstein}

For fixed $(q,r,a)$ write the polynomial as
$P(\xi)=D_0+D_1\xi+D_2\xi^2+D_3\xi^3$.  On a subinterval
$[l,h]$, substitute $\xi=l+(h-l)t$, $0\le t\le1$, and express $P$
in the basis $\binom3j t^j(1-t)^{3-j}$.  The multipliers of $D_1,D_2,D_3$
in the four coefficients are
\begin{equation}
 \begin{array}{c|ccc}
 j& D_1&D_2&D_3\\ \hline
 0&l&l^2&l^3\\
 1&(2l+h)/3&l(l+2h)/3&l^2h\\
 2&(l+2h)/3&h(h+2l)/3&lh^2\\
 3&h&h^2&h^3
 \end{array}
 \label{int:bernstein-rows}
\end{equation}
and the multiplier of $D_0$ is always one.  Expanding the four
basis polynomials verifies these identities.  All basis polynomials
are nonnegative and their sum is one; strictly positive coefficients
therefore imply $P(\xi)>0$ on the whole closed subinterval.

The function \texttt{rows} uses successive intervals with endpoints
\begin{equation}
 \frac1{1000}(30,50,75,100,150,200,300,400,500,650,800,900,1000).
 \label{int:xi-grid}
\end{equation}
These twelve intervals cover $[3/100,1]$ without gaps.  For each it
stores the first three rows of~\eqref{int:bernstein-rows}; the fourth
row equals the first row for the next interval.  It finally adds
$(1,1,1)$.  Thus the $37$ stored rows account for all $48$ coefficients
of the twelve cubic representations, including both extreme
endpoints.  The row multipliers are themselves evaluated by rational
interval arithmetic and have zero gradients with respect to
$(q,r,a)$.

Let $F$ denote one of these Bernstein coefficients on a box
$B=\prod_{i=1}^3[l_i/S,h_i/S]$.  The interval calculation in
\texttt{dot} encloses $F$ and all its first derivatives.  A positive
lower endpoint immediately accepts this coefficient.  Otherwise
define the grid midpoint and coordinate radii by
\[
 c_i=\frac{\lfloor(l_i+h_i)/2\rfloor}{S},\qquad
 \rho_i=\frac{h_i-\lfloor(l_i+h_i)/2\rfloor}{S}.
\]
Nonnegative endpoints ensure that C++ integer division computes this
floor; $\rho_i$ is the larger of the two distances to the endpoints.
If $M_i$ is the maximum absolute endpoint of the derivative enclosure,
in real units, then for every $z\in B$ the segment from $c$ to $z$
lies in $B$ and
\begin{equation}
 F(z)\ge F(c)-\sum_{i=1}^3M_i\rho_i.
 \label{int:mean-value}
\end{equation}
This follows by integrating the gradient along that segment, also
when a coordinate endpoint is zero.  The code encloses $F(c)$ anew,
rounds each product $M_i\rho_i$ upwards, and adds their integer
upper endpoints with overflow checks.  Its strict comparison
\texttt{value.v.l > total} is therefore sufficient for $F>0$
throughout $B$.  A box is accepted only when every stored row passes
one of these two tests.

\subsection{Subdivision, coverage, and observed completion}
\label{int:coverage}

For a band $[a_-,a_+]$, the initial exact boxes before outward
rounding are
\begin{equation}
 \left[\frac i{24},\frac{i+1}{24}\right]
 \times\left[\frac j{24},\frac{j+1}{24}\right]
 \times\left[a_-+\frac{k}{16}(a_+-a_-),
              a_-+\frac{k+1}{16}(a_+-a_-)\right],
 \label{int:initial-boxes}
\end{equation}
where $0\le i<24$, $0\le j\le i$, and $0\le k<16$.
Every $(q,r)$ with $0\le r\le q\le1$ belongs to one of the displayed
$q,r$ rectangles, including all grid boundaries.  There are
$16\sum_{i=0}^{23}(i+1)=4800$ initial boxes per run.  Their endpoints
are rounded outwards, so none of the required domain is lost.
The extra points included by rounding are harmless, as are portions
of diagonal rectangles with $r>q$: the same functions and derivative
bounds are valid there.  The $a$ enlargement and all analytic domains
were checked above.  Adjacent bands may use different parameters,
since each certifies its own entire closed band.

On entry, \texttt{visit} skips a box only if its lower $r$ endpoint
exceeds its upper $q$ endpoint.  Such a box contains no point with
$r\le q$.  If any coefficient fails both acceptance tests, the code
chooses a coordinate with maximal recorded derivative-radius
contribution and cuts at its grid midpoint.  The choice affects
efficiency only.  The two closed children include the midpoint and
their union is the parent.  If the chosen contribution is zero, the
depth exceeds $80$, or the chosen grid radius is at most one, an
assertion fails rather than accepting an unresolved box.

Induction back up any normally completed recursion tree now proves
that every required point in its root box has all Bernstein
coefficients strictly positive.  The leaves either pass the sound
acceptance tests or contain no required point.  Together with
\eqref{int:initial-boxes}, \eqref{int:xi-grid}, and the sorted
parametrization, normal completion of all three runs proves the
pointwise inequality on the full stated domain.  This includes the
uniform message, boundary probability vectors, $a=0$, all band
endpoints, and all $\xi$ endpoints.  In particular there is no
exception at the uniform message: there $P(\xi)=\delta\xi^2+
U\delta\xi^3>0$.

\label{int:execution}
The unchanged source was replayed with GCC 13.3.0 on x86\_64 Linux,
using assertions enabled.  From the package root the reproduction
commands are
\begin{verbatim}
g++ -std=c++17 -O3 -UNDEBUG verification/interval.cpp \
  -o /tmp/appendix_b
/tmp/appendix_b 0
/tmp/appendix_b 1
/tmp/appendix_b 2
\end{verbatim}
Only the documented arguments $0,1,2$ are part of this protocol;
the source does not validate other argument values.  Compiling with
assertions disabled is not this certificate.  All three runs exited
normally with status zero and the following visit counts:
\[
 \begin{array}{c|c|r}
 \text{argument}&a\text{ interval}&\text{calls to \texttt{visit}}\\ \hline
 0&[0,190/1000]&414080\\
 1&[190/1000,448/1000]&756096\\
 2&[448/1000,477/1000]&1004426
 \end{array}
\]
The counts include internal subdivision nodes and boxes skipped by
the diagonal test. All three executions were also repeated on the
unchanged source with \texttt{-O2 -UNDEBUG -fsanitize=undefined}
and \texttt{-fno-sanitize-recover=undefined}. They completed with the
same counts and no sanitizer diagnostics. This supplies the separate
check of the unchecked \texttt{visits++} counter for these prescribed
executions. The counts identify the finite completed executions;
their correctness implication is the enclosure and coverage argument
above, not a sample-based inference.  The source SHA-256 is
\begin{verbatim}
76d1443b1194597ee77749e8030c825f2b9aae46b5c172d9eccf0745078c28e3
\end{verbatim}
